% Część III — interfejsy backendu: REST API, protokoły czasu rzeczywistego, serwer MCP
% Źródła: docs/backend/api-reference.md, docs/backend/realtime.md, backend/mcp_server/README.md,
% README.md (sekcja MCP), zweryfikowane z kodem backend/ (routers/, ws/, serializers.py, config.py).

\chapter{REST API}
\label{ch:api-rest}

Rozdział jest kompletną referencją interfejsu REST platformy PlayLink. Serwer to aplikacja Python
FastAPI uruchamiana przez \code{uv}. Odpowiedzi są w formacie JSON, o ile nie zaznaczono inaczej.
Interaktywna specyfikacja OpenAPI jest dostępna pod adresem \code{/docs} działającego serwera.
Kod endpointów mieści się w modułach \code{backend/routers/} (\code{auth}, \code{users}, \code{rooms},
\code{events}, \code{games}, \code{chat}) oraz w \code{backend/main.py} (sondy zdrowia).
Model danych, koncepcję uwierzytelniania i zmienne konfiguracyjne opisuje część o architekturze;
tutaj podano je tylko w zakresie niezbędnym do korzystania z API.

\section{Informacje ogólne}
\label{sec:api-ogolne}

\subsection{Adresy bazowe}

\begin{tabularx}{\textwidth}{lY}
\toprule
Środowisko & Adres \\
\midrule
Lokalne (rozwój) & \code{http://localhost:8000} \\
Produkcja & \code{https://playlink.bartek.monster} \\
\bottomrule
\end{tabularx}

\subsection{Uwierzytelnianie}
\label{sec:api-auth}

Endpointy wymagające uwierzytelnienia oczekują nagłówka HTTP \code{Authorization} w schemacie
OAuth2 \emph{Password Bearer}:

\begin{lstlisting}
Authorization: Bearer <jwt>
\end{lstlisting}

Token JWT uzyskuje się z \code{POST /auth/verify} (sekcja o grupie Auth). Jest to standardowy
JWT HS256 o konfigurowalnym czasie życia (domyślnie 60 minut, zmienna \code{JWT_EXPIRATION_MINUTES}).
Zależność FastAPI to \code{OAuth2PasswordBearer(tokenUrl="auth/verify", auto_error=False)}:
brak lub nieprawidłowy token daje \code{401} bez przekierowania. Endpointy oznaczone w tabeli jako
\emph{Admin} wymagają dodatkowo, by adres z tokenu znajdował się na liście \code{ADMIN_ADDRESSES};
lista jest autorytatywna, więc sfałszowany claim \code{is_admin} w JWT nie daje dostępu.

\subsection{Ograniczanie liczby żądań}

Każdy endpoint REST (poza sondami \code{/healthz} i \code{/readyz}) jest chroniony przez
\code{slowapi} z limitem domyślnym ze zmiennej \code{DEFAULT_RATE_LIMIT} (domyślnie
\code{"10/minute"}). Przekroczenie limitu daje odpowiedź \code{429 Too Many Requests}
z treścią w rodzaju \code{{"detail": "Rate limit exceeded: 10 per 1 minute"}}
(lub z wartością skonfigurowaną).

\subsection{Błędy globalne}

\begin{tabularx}{\textwidth}{lY}
\toprule
Status & Warunek \\
\midrule
\code{422 Unprocessable Entity} & Treść żądania nie przechodzi walidacji Pydantic (brakujące pola, błędy typów, naruszenie ograniczeń); zwracany automatycznie przez FastAPI. \\
\code{429 Too Many Requests} & Przekroczony limit żądań. \\
\bottomrule
\end{tabularx}

\section{Przegląd endpointów}
\label{sec:api-przeglad}

Kolumna \emph{Autoryzacja}: „---” (publiczny), „JWT” (wymagany token), „Admin” (token i adres na
liście administratorów).

\begin{longtable}{@{}l p{5.6cm} l p{3.4cm}@{}}
\toprule
Metoda & Ścieżka & Autoryzacja & Opis \\
\midrule
\endhead
\code{GET} & \code{/} & --- & Zgodna wstecznie sonda zdrowia \\
\code{GET} & \code{/healthz} & --- & Sonda żywotności z regionem i wersją (bez limitu żądań) \\
\code{GET} & \code{/readyz} & --- & Sonda gotowości z kontrolą bazy (bez limitu żądań) \\
\code{POST} & \code{/auth/request-nonce} & --- & Jednorazowy nonce uwierzytelniający \\
\code{POST} & \code{/auth/verify} & --- & Weryfikacja podpisu i wydanie JWT \\
\code{GET} & \code{/users/me} & JWT & Profil zalogowanego użytkownika \\
\code{PATCH} & \code{/users/me} & JWT & Zmiana nazwy użytkownika \\
\code{GET} & \code{/rooms} & --- & Lista aktywnych pokoi \\
\code{GET} & \code{/lobby-locations} & --- & Lista obsługiwanych lokalizacji lobby \\
\code{GET} & \code{/rooms/{room_name}} & --- & Szczegóły jednego pokoju \\
\code{GET} & \code{/games} & --- & Lista kategorii gier \\
\code{POST} & \code{/rooms} & JWT & Utworzenie pokoju \\
\code{POST} & \code{/rooms/{room_name}/join} & JWT & Dołączenie do pokoju \\
\code{POST} & \code{/rooms/{room_name}/leave} & JWT & Opuszczenie pokoju \\
\code{POST} & \code{/rooms/{room_name}/members/}\newline\code{{member_address}/kick} & Admin & Usunięcie członka z pokoju \\
\code{GET} & \code{/rooms/{room_name}/event} & --- & Zaplanowane wydarzenie pokoju \\
\code{PUT} & \code{/rooms/{room_name}/event} & JWT & Zaplanowanie lub zmiana wydarzenia \\
\code{DELETE} & \code{/rooms/{room_name}/event} & JWT & Odwołanie wydarzenia \\
\code{PUT} & \code{/rooms/{room_name}/event/rsvp} & JWT & Ustawienie RSVP wywołującego \\
\code{POST} & \code{/games} & Admin & Dodanie kategorii gry \\
\code{DELETE} & \code{/games/{name}} & Admin & Usunięcie kategorii gry \\
\code{DELETE} & \code{/rooms/{room_name}} & Admin & Zamknięcie pokoju \\
\bottomrule
\end{longtable}

Endpointy WebSocket (\code{/ws/rooms}, \code{/ws/rooms/{room_name}/chat}) opisuje kolejny rozdział.
Poza tabelą istnieje jeszcze \code{GET /metrics} (metryki Prometheus), wyłączony ze schematu OpenAPI:
jest odczytywany wyłącznie w sieci Compose, a żądania przychodzące przez Cloudflare
(nagłówek \code{cf-connecting-ip}) dostają \code{404}. Szczegóły w części o obserwowalności.

\section{Zdrowie (Health)}
\label{sec:api-health}

\subsection{\texttt{GET /}}

Podstawowa sonda zdrowia, bez uwierzytelniania.

\begin{lstlisting}[language=Python]
200 OK
{ "status": "ok", "service": "playlink-auth" }
\end{lstlisting}

\subsection{\texttt{GET /healthz}}

Tania sonda żywotności używana przez selektor regionu i healthchecki kontenerów. Nie dotyka bazy
danych i celowo \emph{nie} podlega limitowi żądań. Odpowiedź ma nagłówek \code{Cache-Control: no-store}.
Pola \code{region}, \code{version} i \code{commit} pochodzą ze zmiennych \code{REGION},
\code{APP_VERSION} i \code{GIT_SHA}.

\begin{lstlisting}[language=Python]
200 OK
{ "status": "ok", "region": "eu", "version": "v0.1.0", "commit": "abc1234" }
\end{lstlisting}

\subsection{\texttt{GET /readyz}}

Sonda gotowości wykonująca \code{SELECT 1} w PostgreSQL. Gdy baza jest niedostępna, zwraca
\code{503} (\code{{"detail": "database unavailable"}}). Używana przez healthcheck Compose;
sonda regionu w przeglądarce powinna korzystać z \code{/healthz}.

\begin{lstlisting}[language=Python]
200 OK
{ "status": "ready", "region": "eu" }
\end{lstlisting}

\section{Auth}
\label{sec:api-auth-endpointy}

\subsection{\texttt{POST /auth/request-nonce}}

Żąda jednorazowego nonce'a dla adresu tożsamości Ethereum. Nonce jest przed zapisem haszowany
(SHA-256); wartość jawna jest zwracana tylko raz i musi zostać podpisana EIP-191 \code{personal_sign}
w kroku weryfikacji. Poprzednie niewykorzystane nonce'y tego adresu są automatycznie unieważniane.
Jeżeli użytkownik o danym adresie nie istnieje, rekord jest tworzony (upsert). Nonce wygasa po
\code{NONCE_EXPIRATION_MINUTES} (domyślnie 5).

\textbf{Parametr zapytania}

\begin{tabularx}{\textwidth}{llY}
\toprule
Pole & Typ & Ograniczenia \\
\midrule
\code{address} & \code{string} & Poprawny adres Ethereum (hex, prefiks \code{0x}, 42 znaki); konwertowany do sumy kontrolnej EIP-55. \\
\bottomrule
\end{tabularx}

\begin{lstlisting}[language=Python]
200 OK
{ "nonce": "550e8400-e29b-41d4-a716-446655440000" }
\end{lstlisting}

\textbf{Błędy:} \code{400 Bad Request} --- \code{address} nie ma poprawnego formatu
(\code{{"detail": "Invalid identity address format"}}).

\subsection{\texttt{POST /auth/verify}}

Weryfikuje podpis EIP-191 \code{personal_sign} złożony pod nonce'em i wydaje JWT. Podpisywany
komunikat ma postać:

\begin{lstlisting}
Sign in to Playlink\nNonce: {nonce}
\end{lstlisting}

Po sukcesie nonce zostaje oznaczony jako użyty, pole \code{last_login} użytkownika jest
aktualizowane i zwracany jest token.

\textbf{Treść żądania (\code{VerifyRequest})}

\begin{tabularx}{\textwidth}{llY}
\toprule
Pole & Typ & Ograniczenia \\
\midrule
\code{address} & \code{string} & Poprawny adres Ethereum (zgodny z sumą kontrolną EIP-55). \\
\code{nonce} & \code{string} & Jawny nonce zwrócony przez \code{POST /auth/request-nonce}. \\
\code{signature} & \code{string} & Podpis EIP-191 w hex (prefiks \code{0x}, 132 znaki) komunikatu powyżej, złożony kluczem prywatnym odpowiadającym \code{address}. \\
\bottomrule
\end{tabularx}

\begin{lstlisting}[language=Python]
200 OK
{ "token": "eyJhbGciOiJIUzI1NiIs...", "username": "alice_42", "is_admin": false }
\end{lstlisting}

\begin{tabularx}{\textwidth}{llY}
\toprule
Pole & Typ & Opis \\
\midrule
\code{token} & \code{string} & JWT HS256. Czas życia: \code{JWT_EXPIRATION_MINUTES} (domyślnie 60 min). Claimy: \code{sub} (adres w sumie kontrolnej), \code{username}, \code{is_admin}, \code{iat}, \code{exp}, \code{iss} (\code{"playlink-auth"}). \\
\code{username} & \code{string} & Nazwa wyświetlana; generowana automatycznie, jeśli nie ustawiona (format \code{user_<8 znaków hex>}). \\
\code{is_admin} & \code{boolean} & Czy adres jest na liście \code{ADMIN_ADDRESSES}. Klient nie może opierać się na claimach JWT, by omijać tę kontrolę na endpointach administracyjnych. \\
\bottomrule
\end{tabularx}

\textbf{Błędy}

\begin{tabularx}{\textwidth}{lY}
\toprule
Status & Warunek \\
\midrule
\code{400} & Nieprawidłowy format \code{address} (\code{"Invalid identity address format"}). \\
\code{401} & Nonce nie istnieje, został już użyty lub wygasł (\code{"Invalid or expired challenge"}). \\
\code{401} & Nieprawidłowy format podpisu --- odzyskanie adresu się nie powiodło (\code{"Invalid signature format"}). \\
\code{401} & Odzyskany adres nie zgadza się z \code{address} (\code{"Identity verification failed"}). \\
\bottomrule
\end{tabularx}

\section{Użytkownicy}
\label{sec:api-users}

\subsection{\texttt{GET /users/me}}

Zwraca profil uwierzytelnionego użytkownika (wymagany JWT).

\begin{lstlisting}[language=Python]
200 OK
{
  "id": 1,
  "identity_address": "0x742d35Cc6634C0532925a3b844Bc9e7595f2bD18",
  "username": "alice_42",
  "created_at": "2026-06-18T12:00:00Z",
  "last_login": "2026-06-18T14:30:00Z",
  "is_admin": false
}
\end{lstlisting}

\begin{tabularx}{\textwidth}{llY}
\toprule
Pole & Typ & Opis \\
\midrule
\code{id} & \code{integer} & Wewnętrzny identyfikator użytkownika. \\
\code{identity_address} & \code{string} & Adres Ethereum (suma kontrolna EIP-55). \\
\code{username} & \code{string | null} & Nazwa wyświetlana lub \code{null}, jeśli nie ustawiona. \\
\code{created_at} & \code{string} & Znacznik czasu ISO 8601 z sufiksem \code{Z}. \\
\code{last_login} & \code{string | null} & Czas ostatniego udanego \code{POST /auth/verify}. \\
\code{is_admin} & \code{boolean} & Czy adres jest na liście \code{ADMIN_ADDRESSES}. \\
\bottomrule
\end{tabularx}

\textbf{Błędy:} \code{401} --- brak lub nieprawidłowy JWT (\code{"Not authenticated"},
\code{"Token expired"}, \code{"Invalid token"}); \code{404} --- brak rekordu użytkownika
(\code{"User not found"}; po poprawnym uwierzytelnieniu nie powinno się zdarzyć).

\subsection{\texttt{PATCH /users/me}}

Zmienia nazwę użytkownika (wymagany JWT). Treść żądania (\code{UpdateUserRequest}):
pole \code{username} (\code{string}), 3--20 znaków; dozwolone litery \code{a-zA-Z}, cyfry,
podkreślenie \code{_} i łącznik \code{-}. Nazwa nie może pasować do żadnego wpisu listy wulgaryzmów
(LDNOOBW, angielska), sprawdzanej dla całego ciągu oraz dla każdego tokenu rozdzielonego znakami
\code{_}/\code{-}. Ustawienie nazwy identycznej z obecną zwraca profil bez zmian.

Odpowiedź \code{200 OK} ma taki sam kształt jak \code{GET /users/me}.

\begin{tabularx}{\textwidth}{lY}
\toprule
Status & Warunek \\
\midrule
\code{400} & Błędny format nazwy (\code{"Username must be 3-20 characters: letters, numbers, _ or -."}). \\
\code{400} & Nazwa zawiera wulgaryzmy (\code{"Username contains inappropriate language."}). \\
\code{401} & Brak lub nieprawidłowy JWT. \\
\code{404} & Brak rekordu użytkownika. \\
\code{409} & Nazwa zajęta przez innego użytkownika (\code{"Username already taken."}). \\
\bottomrule
\end{tabularx}

\section{Pokoje}
\label{sec:api-rooms}

\subsection{\texttt{GET /rooms}}

Lista aktywnych (niewygasłych) pokoi; bez uwierzytelniania. Endpoint zwraca surowe obiekty modelu ORM
\code{Room} jako tablicę JSON.

\begin{lstlisting}[language=Python]
200 OK
[
  {
    "name": "dm-arena",
    "game": "Quake III Arena",
    "lobby_location": "eu-central",
    "players_max": 8,
    "description": "Casual deathmatch",
    "communicator_link": "https://discord.gg/example",
    "requirements": "Microphone preferred",
    "created_by": "0x742d35Cc6634C0532925a3b844Bc9e7595f2bD18",
    "created_at": "2026-06-18T12:00:00Z",
    "expires_at": "2026-06-18T13:00:00Z",
    "members": []
  }
]
\end{lstlisting}

Bogatszy widok (z \code{member_addresses}, \code{players_active} i zagnieżdżonym \code{event}) daje
\code{GET /rooms/{room_name}}.

\begin{tabularx}{\textwidth}{llY}
\toprule
Pole & Typ & Opis \\
\midrule
\code{name} & \code{string} & Unikalna nazwa pokoju. \\
\code{game} & \code{string} & Nazwa kategorii gry. \\
\code{lobby_location} & \code{string} & Kod lokalizacji lobby (np. \code{"eu-central"}). \\
\code{players_max} & \code{integer} & Maksymalna liczba członków. \\
\code{description} & \code{string | null} & Opis (maks. 500 znaków). \\
\code{communicator_link} & \code{string | null} & Link do komunikatora głosowego/tekstowego (maks. 500 znaków). \\
\code{requirements} & \code{string | null} & Wymagania wobec graczy (maks. 1000 znaków). \\
\code{created_by} & \code{string} & Adres tożsamości twórcy pokoju. \\
\code{created_at} & \code{string} & Znacznik czasu ISO 8601. \\
\code{expires_at} & \code{string} & Znacznik czasu ISO 8601. Domyślnie pokój wygasa 60 minut po utworzeniu (automatycznie przedłużane po zaplanowaniu wydarzenia). \\
\code{members} & \code{array} & Lista obiektów \code{User}; w praktyce pusta tablica --- skład pokoju zwraca \code{GET /rooms/{room_name}}. \\
\bottomrule
\end{tabularx}

\subsection{\texttt{GET /rooms/\{room\_name\}}}

Pełne szczegóły pokoju, w tym lista członków i zaplanowane wydarzenie; bez uwierzytelniania.
Parametr ścieżki \code{room_name} to unikalna nazwa pokoju.

\begin{lstlisting}[language=Python]
200 OK
{
  "name": "dm-arena",
  "game": "Quake III Arena",
  "lobby_location": "eu-central",
  "players_max": 8,
  "players_active": 2,
  "member_addresses": [
    "0x742d35Cc6634C0532925a3b844Bc9e7595f2bD18",
    "0x1234567890abcdef1234567890abcdef12345678"
  ],
  "members": [
    { "address": "0x742d35Cc6634C0532925a3b844Bc9e7595f2bD18",
      "username": "alice_42", "is_admin": false }
  ],
  "description": "Casual deathmatch",
  "communicator_link": "https://discord.gg/example",
  "requirements": "Microphone preferred",
  "created_by": "0x742d35Cc6634C0532925a3b844Bc9e7595f2bD18",
  "expires_at": "2026-06-18T13:00:00Z",
  "event": null
}
\end{lstlisting}

\begin{tabularx}{\textwidth}{llY}
\toprule
Pole & Typ & Opis \\
\midrule
\code{name}, \code{game}, \code{lobby_location}, \code{players_max} & & Jak w \code{GET /rooms}. \\
\code{players_active} & \code{integer} & Bieżąca liczba członków. \\
\code{member_addresses} & \code{array[string]} & Adresy tożsamości wszystkich członków. \\
\code{members} & \code{array[object]} & Szczegółowa lista członków: \code{{ address, username, is_admin }}. \\
\code{description}, \code{communicator_link}, \code{requirements} & \code{string | null} & Opis, link do komunikatora, wymagania. \\
\code{created_by} & \code{string} & Adres twórcy. \\
\code{expires_at} & \code{string} & Czas wygaśnięcia (ISO 8601). \\
\code{event} & \code{object | null} & Zaplanowane wydarzenie lub \code{null}; kształt jak w odpowiedzi \code{GET /rooms/{room_name}/event}. \\
\bottomrule
\end{tabularx}

\textbf{Błędy:} \code{404} --- pokój nie istnieje (\code{"Room not found"}).

\subsection{\texttt{POST /rooms}}

Tworzy nowy pokój; twórca jest automatycznie dodawany jako członek. Wymagany JWT.

\textbf{Efekty uboczne.} Rozgłasza zaktualizowaną listę pokoi przez \code{/ws/rooms} do wszystkich
słuchaczy kanału globalnego. Jeśli nazwa gry jest nowa, rekord \code{Game} powstaje automatycznie
(\code{_ensure_game}, na końcu kolejności sortowania). Przed sprawdzeniem unikalności nazwy
endpoint usuwa wygasłe pokoje (wraz z ich wiadomościami).

\textbf{Limit.} Użytkownik niebędący administratorem może mieć najwyżej \textbf{3 aktywne pokoje}
(liczone są pokoje, w których \code{created_by} równa się adresowi wywołującego).

\textbf{Treść żądania (\code{CreateRoomRequest})}

\begin{tabularx}{\textwidth}{llY}
\toprule
Pole & Typ & Ograniczenia \\
\midrule
\code{name} & \code{string} & Unikalna nazwa pokoju; nie może już istnieć. \\
\code{game} & \code{string} & Nazwa kategorii gry, 1--100 znaków; nieistniejąca gra jest tworzona automatycznie. \\
\code{lobby_location} & \code{string} & Jeden z 10 obsługiwanych kodów (patrz \code{GET /lobby-locations}); 1--50 znaków. \\
\code{players_max} & \code{integer} & Maksymalna liczba członków (API nie narzuca poza typem \code{integer} żadnego minimum ani maksimum). \\
\code{description} & \code{string | null} & Maks. 500 znaków. \\
\code{communicator_link} & \code{string | null} & Poprawny URL (Pydantic \code{HttpUrl}), jeśli podany; zapisywany jako tekst. \\
\code{requirements} & \code{string | null} & Maks. 1000 znaków. \\
\bottomrule
\end{tabularx}

\begin{lstlisting}[language=Python]
201 Created
{ "status": "created", "room": "dm-arena" }
\end{lstlisting}

\begin{tabularx}{\textwidth}{lY}
\toprule
Status & Warunek \\
\midrule
\code{400} & Osiągnięty limit 3 pokoi (\code{"You can create a maximum of 3 rooms."}). \\
\code{400} & Pusta nazwa gry po przycięciu (\code{"Game name is required"}). \\
\code{400} & Nieznany kod lokalizacji (\code{"Unsupported lobby location"}). \\
\code{401} & Brak lub nieprawidłowy JWT. \\
\code{409} & Nazwa pokoju zajęta (\code{"Room name already taken"}). \\
\code{422} & Błąd walidacji Pydantic (np. \code{players_max} nie jest liczbą, \code{communicator_link} nie jest URL-em, \code{description} dłuższy niż 500 znaków). \\
\bottomrule
\end{tabularx}

\subsection{\texttt{POST /rooms/\{room\_name\}/join}}

Dołącza wywołującego do pokoju (wymagany JWT).

\textbf{Efekty uboczne:} rozgłoszenie listy pokoi przez \code{/ws/rooms} oraz ramka \code{roster_update}
na kanał czatu pokoju.

\begin{lstlisting}[language=Python]
200 OK
{ "status": "joined", "room": "dm-arena" }
\end{lstlisting}

\begin{tabularx}{\textwidth}{lY}
\toprule
Status & Warunek \\
\midrule
\code{400} & Wywołujący jest już członkiem (\code{"You are already in this room"}). \\
\code{400} & Pokój pełny --- liczba członków $\geq$ \code{players_max} (\code{"Room is full"}). \\
\code{401} & Brak lub nieprawidłowy JWT. \\
\code{404} & Pokój nie istnieje (\code{"Room not found"}) lub brak rekordu użytkownika (\code{"User not found"}). \\
\bottomrule
\end{tabularx}

\subsection{\texttt{POST /rooms/\{room\_name\}/leave}}

Opuszcza pokój (wymagany JWT). RSVP wywołującego na wydarzenie pokoju (jeśli istnieje) jest
automatycznie usuwane.

\textbf{Efekty uboczne:} rozgłoszenie listy pokoi przez \code{/ws/rooms}; ramka \code{roster_update}
na kanał czatu; jeżeli opuszczający miał RSVP --- także ramka \code{event_update} z nowym stanem wydarzenia.

\begin{lstlisting}[language=Python]
200 OK
{ "status": "left", "room": "dm-arena" }
\end{lstlisting}

\textbf{Błędy:} \code{400} --- wywołujący nie jest członkiem (\code{"You are not in this room"});
\code{401} --- brak/nieprawidłowy JWT; \code{404} --- pokój lub użytkownik nie istnieje
(\code{"Room not found"}, \code{"User not found"}).

\subsection{\texttt{POST /rooms/\{room\_name\}/members/\{member\_address\}/kick}}

Usuwa z pokoju członka niebędącego administratorem. Wyłącznie administratorzy (adresy z
\code{ADMIN_ADDRESSES}). Parametr \code{member_address} to adres tożsamości usuwanego członka;
wyszukiwanie nie rozróżnia wielkości liter.

\textbf{Efekty uboczne}
\begin{itemize}[nosep]
  \item Jeśli usunięty był twórcą pokoju, własność przechodzi na kopiącego administratora (zmiana \code{room.created_by}).
  \item RSVP usuniętego członka (jeśli istnieje) jest kasowane.
  \item W czacie pokoju zapisywana jest wiadomość systemowa: \code{"{username} was removed from the room by an administrator."}
  \item Na kanał czatu pokoju trafiają ramki (w tej kolejności): \code{message}, \code{roster_update}, \code{event_update} (jeśli istnieje wydarzenie) i \code{member_kicked}.
  \item Rozgłoszenie zaktualizowanej listy pokoi przez \code{/ws/rooms}.
  \item Aktywne połączenia czatu usuniętego członka są zamykane kodem \code{4409}.
\end{itemize}

\begin{lstlisting}[language=Python]
200 OK
{
  "status": "kicked",
  "member_address": "0x1234567890abcdef1234567890abcdef12345678",
  "member_username": "bob_99",
  "created_by": "0x742d35Cc6634C0532925a3b844Bc9e7595f2bD18",
  "ownership_transferred": false
}
\end{lstlisting}

\begin{tabularx}{\textwidth}{llY}
\toprule
Pole & Typ & Opis \\
\midrule
\code{status} & \code{string} & \code{"kicked"}. \\
\code{member_address} & \code{string} & Adres usuniętego członka. \\
\code{member_username} & \code{string} & Nazwa usuniętego członka. \\
\code{created_by} & \code{string} & Twórca pokoju po operacji (mógł się zmienić). \\
\code{ownership_transferred} & \code{boolean} & Czy własność pokoju przeszła na administratora. \\
\bottomrule
\end{tabularx}

\begin{tabularx}{\textwidth}{lY}
\toprule
Status & Warunek \\
\midrule
\code{401} & Brak lub nieprawidłowy JWT. \\
\code{403} & Wywołujący nie jest administratorem (\code{"Admin privileges required"}). \\
\code{403} & Cel jest administratorem (\code{"Administrators cannot be kicked"}). \\
\code{404} & Pokój nie istnieje (\code{"Room not found"}) lub użytkownik nie istnieje (\code{"User not found"}). \\
\code{409} & Cel nie jest członkiem pokoju (\code{"User is not a room member"}). \\
\bottomrule
\end{tabularx}

\subsection{\texttt{DELETE /rooms/\{room\_name\}}}

Zamyka (usuwa) pokój; wyłącznie administrator. Kaskadowo usuwa wiadomości czatu, zaplanowane
wydarzenie i wszystkie RSVP.

\textbf{Efekty uboczne:} ramka \code{{"type": "room_closed", "room": "{room_name}"}} na kanał czatu
pokoju (klienci są przekierowywani) oraz rozgłoszenie listy pokoi przez \code{/ws/rooms}.

\begin{lstlisting}[language=Python]
200 OK
{ "status": "closed", "room": "dm-arena" }
\end{lstlisting}

\textbf{Błędy:} \code{401} --- brak/nieprawidłowy JWT; \code{403} --- brak uprawnień administratora
(\code{"Admin privileges required"}); \code{404} --- pokój nie istnieje (\code{"Room not found"}).

\section{Wydarzenia pokoju i RSVP}
\label{sec:api-events}

Każdy pokój ma co najwyżej jedno zaplanowane wydarzenie. Przeglądać harmonogram może każdy, ale
RSVP mogą ustawiać tylko członkowie, a planować i odwoływać --- tylko twórca pokoju.

\subsection{\texttt{GET /rooms/\{room\_name\}/event}}

Zwraca zaplanowane wydarzenie razem z listą RSVP; endpoint publiczny.

\begin{lstlisting}[language=Python]
200 OK
{
  "starts_at": "2026-06-20T18:00:00Z",
  "ends_at": "2026-06-20T20:00:00Z",
  "created_by": "0x742d35Cc6634C0532925a3b844Bc9e7595f2bD18",
  "created_at": "2026-06-18T12:00:00Z",
  "updated_at": "2026-06-18T12:00:00Z",
  "rsvps": [
    { "address": "0x742d35Cc6634C0532925a3b844Bc9e7595f2bD18",
      "username": "alice_42", "status": "present",
      "updated_at": "2026-06-18T12:05:00Z" }
  ]
}
\end{lstlisting}

\begin{tabularx}{\textwidth}{llY}
\toprule
Pole & Typ & Opis \\
\midrule
\code{starts_at}, \code{ends_at} & \code{string} & Początek i koniec wydarzenia (ISO 8601). \\
\code{created_by} & \code{string} & Adres twórcy pokoju, który zaplanował wydarzenie. \\
\code{created_at}, \code{updated_at} & \code{string} & Znaczniki utworzenia i ostatniej zmiany. \\
\code{rsvps} & \code{array[object]} & Lista RSVP: \code{{ address, username, status ("present"|"absent"|"maybe"), updated_at }}. \\
\bottomrule
\end{tabularx}

\textbf{Błędy:} \code{404} --- pokój nie istnieje (\code{"Room not found"}) lub brak wydarzenia
(\code{"No event scheduled"}).

\subsection{\texttt{PUT /rooms/\{room\_name\}/event}}

Tworzy lub zastępuje wydarzenie pokoju. Tylko twórca pokoju (adres z JWT musi równać się
\code{room.created_by}, bez rozróżniania wielkości liter). \code{expires_at} pokoju jest automatycznie
przedłużane do \code{ends_at} + 30 minut (okres karencji), o ile nowy termin jest późniejszy.
Jeśli okno czasowe istniejącego wydarzenia ulega zmianie, wszystkie dotychczasowe RSVP są usuwane
--- członkowie muszą potwierdzić obecność ponownie.

\textbf{Efekty uboczne:} ramka \code{event_update} na kanał czatu pokoju; gdy zmienił się czas
wygaśnięcia pokoju --- także rozgłoszenie listy pokoi przez \code{/ws/rooms}.

\textbf{Treść żądania (\code{ScheduleEventRequest})}

\begin{tabularx}{\textwidth}{llY}
\toprule
Pole & Typ & Ograniczenia \\
\midrule
\code{starts_at} & \code{datetime} & Musi leżeć w przyszłości (UTC); wartość bez strefy jest traktowana jako UTC. \\
\code{ends_at} & \code{datetime} & Musi być późniejsze niż \code{starts_at}; wartość bez strefy jest traktowana jako UTC. \\
\bottomrule
\end{tabularx}

Odpowiedź \code{200 OK} ma taki sam kształt jak \code{GET /rooms/{room_name}/event}.

\begin{tabularx}{\textwidth}{lY}
\toprule
Status & Warunek \\
\midrule
\code{400} & \code{starts_at} nie leży w przyszłości (\code{"starts_at must be in the future"}). \\
\code{400} & \code{ends_at} nie jest późniejsze niż \code{starts_at} (\code{"ends_at must be after starts_at"}). \\
\code{401} & Brak lub nieprawidłowy JWT. \\
\code{403} & Wywołujący nie jest twórcą (\code{"Only the room creator can schedule an event"}). \\
\code{404} & Pokój nie istnieje (\code{"Room not found"}). \\
\bottomrule
\end{tabularx}

\subsection{\texttt{DELETE /rooms/\{room\_name\}/event}}

Odwołuje wydarzenie i kasuje wszystkie RSVP; tylko twórca pokoju. \textbf{Efekt uboczny:} ramka
\code{{"type": "event_update", "event": null}} na kanał czatu pokoju.

\begin{lstlisting}[language=Python]
200 OK
{ "status": "cancelled", "room": "dm-arena" }
\end{lstlisting}

\textbf{Błędy:} \code{401} --- brak/nieprawidłowy JWT; \code{403} --- wywołujący nie jest twórcą
(\code{"Only the room creator can cancel the event"}); \code{404} --- pokój nie istnieje lub brak
wydarzenia (\code{"Room not found"}, \code{"No event scheduled"}).

\subsection{\texttt{PUT /rooms/\{room\_name\}/event/rsvp}}

Ustawia (upsert) RSVP wywołującego na wydarzenie pokoju. Tylko aktualni członkowie pokoju; jedno
RSVP na użytkownika na wydarzenie (poprzedni status jest nadpisywany). \textbf{Efekt uboczny:} ramka
\code{rsvp_update} na kanał czatu pokoju.

Treść żądania (\code{SetRsvpRequest}): pole \code{status} (\code{string}) o wartości \code{"present"},
\code{"absent"} lub \code{"maybe"} (wartości enuma \code{RsvpStatus}).

\begin{lstlisting}[language=Python]
200 OK
{
  "address": "0x742d35Cc6634C0532925a3b844Bc9e7595f2bD18",
  "username": "alice_42",
  "status": "present",
  "updated_at": "2026-06-18T12:05:00Z"
}
\end{lstlisting}

\begin{tabularx}{\textwidth}{lY}
\toprule
Status & Warunek \\
\midrule
\code{401} & Brak lub nieprawidłowy JWT. \\
\code{403} & Wywołujący nie jest członkiem (\code{"Only room members can RSVP to the event"}). \\
\code{404} & Pokój, użytkownik lub wydarzenie nie istnieje (\code{"Room not found"}, \code{"User not found"}, \code{"No event scheduled"}). \\
\bottomrule
\end{tabularx}

\section{Lokalizacje lobby}
\label{sec:api-lobby}

\subsection{\texttt{GET /lobby-locations}}

Zwraca listę obsługiwanych lokalizacji lobby z etykietami i przybliżonymi współrzędnymi
geograficznymi oraz kod lokalizacji domyślnej (\code{"eu-central"}). Lista jest zdefiniowana
w \code{backend/config.py} (\code{LOBBY_LOCATIONS}).

\begin{lstlisting}[language=Python]
200 OK
{
  "default": "eu-central",
  "locations": [
    { "code": "na-east",      "label": "North America East",  "lat": 39.0,  "lon": -77.0 },
    { "code": "na-west",      "label": "North America West",  "lat": 37.8,  "lon": -122.4 },
    { "code": "eu-west",      "label": "Europe West",         "lat": 51.5,  "lon": -0.1 },
    { "code": "eu-central",   "label": "Europe Central",      "lat": 50.1,  "lon": 8.7 },
    { "code": "eu-north",     "label": "Europe North",        "lat": 59.3,  "lon": 18.1 },
    { "code": "sa-east",      "label": "South America East",  "lat": -23.6, "lon": -46.6 },
    { "code": "asia-east",    "label": "Asia East",           "lat": 35.7,  "lon": 139.7 },
    { "code": "asia-south",   "label": "Asia South",          "lat": 1.3,   "lon": 103.8 },
    { "code": "oceania",      "label": "Oceania",             "lat": -33.9, "lon": 151.2 },
    { "code": "africa-south", "label": "Africa South",        "lat": -26.2, "lon": 28.0 }
  ]
}
\end{lstlisting}

Pola: \code{default} (\code{string}) --- domyślny kod lokalizacji; \code{locations} (\code{array[object]})
--- obiekty \code{{ code, label, lat, lon }}.

\section{Gry}
\label{sec:api-games}

\subsection{\texttt{GET /games}}

Lista kategorii gier uporządkowana wg \code{sort_order}. Pięć gier domyślnych jest zasiewanych
przy pierwszym uruchomieniu; kolejne dodają i usuwają endpointy administracyjne.

\begin{lstlisting}[language=Python]
200 OK
["Quake III Arena", "Diablo II", "StarCraft", "Half-Life", "Unreal Tournament"]
\end{lstlisting}

\subsection{\texttt{POST /games}}

Dodaje kategorię gry na końcu kolejności sortowania. Wymagany administrator. Treść żądania
(\code{CreateGameRequest}): \code{name} (\code{string}), 1--100 znaków, przycinana po stronie serwera.

\begin{lstlisting}[language=Python]
201 Created
{ "name": "Team Fortress 2", "sort_order": 6 }
\end{lstlisting}

\textbf{Błędy:} \code{400} --- pusta nazwa po przycięciu (\code{"Game name is required"});
\code{401} --- brak/nieprawidłowy JWT; \code{403} --- brak uprawnień administratora
(\code{"Admin privileges required"}); \code{409} --- gra już istnieje (\code{"Game already exists"}).

\subsection{\texttt{DELETE /games/\{name\}}}

Usuwa kategorię gry. Jeśli istnieją pokoje z tą grą, żądanie jest odrzucane kodem \code{409}, chyba
że ustawiono parametr zapytania \code{force=true}; wtedy pokoje są najpierw zamykane (taka sama
kaskada jak w \code{DELETE /rooms/{room_name}}). Wymagany administrator.

\textbf{Efekty uboczne (przy \code{force=true}):} każdy dotknięty pokój jest czyszczony (wiadomości,
wydarzenie, RSVP, członkowie), na kanał czatu każdego z nich trafia ramka \code{room_closed}, a lista
pokoi jest rozgłaszana przez \code{/ws/rooms}.

\begin{tabularx}{\textwidth}{llllY}
\toprule
Parametr & Miejsce & Typ & Domyślnie & Opis \\
\midrule
\code{name} & ścieżka & \code{string} & --- & Nazwa gry. \\
\code{force} & zapytanie & \code{boolean} & \code{false} & Jeśli \code{true}, zamyka wszystkie pokoje z tą grą przed jej usunięciem. \\
\bottomrule
\end{tabularx}

\begin{lstlisting}[language=Python]
200 OK
{ "status": "deleted", "game": "Team Fortress 2", "rooms_closed": ["dm-arena", "ctf-base"] }
\end{lstlisting}

Pole \code{rooms_closed} (\code{array[string]}) to nazwy zamkniętych pokoi (pusta tablica, gdy brak).

\begin{tabularx}{\textwidth}{lY}
\toprule
Status & Warunek \\
\midrule
\code{401} & Brak lub nieprawidłowy JWT. \\
\code{403} & Brak uprawnień administratora (\code{"Admin privileges required"}). \\
\code{404} & Gra nie istnieje (\code{"Game not found"}). \\
\code{409} & Istnieją aktywne pokoje z tą grą, a \code{force} nie jest \code{true} (\code{"There are N active rooms currently playing this game."}). \\
\bottomrule
\end{tabularx}

\section{Odesłania}

Encje (\code{Room}, \code{RoomMember}, \code{Message}, \code{RoomEvent}, \code{RoomEventRsvp},
\code{User}, \code{Nonce}, \code{Game}) opisuje rozdział o modelu danych w części o architekturze;
zmienne środowiskowe (\code{DATABASE_URL}, \code{JWT_SECRET}, \code{JWT_ALGORITHM},
\code{NONCE_EXPIRATION_MINUTES}, \code{JWT_EXPIRATION_MINUTES}, \code{DEFAULT_RATE_LIMIT},
\code{ADMIN_ADDRESSES}, \code{ROOM_CLEANUP_INTERVAL_SECONDS} i inne) --- rozdział o konfiguracji
w tej samej części.

% =====================================================================

\chapter{Komunikacja w czasie rzeczywistym}
\label{ch:realtime}

PlayLink udostępnia dwa endpointy WebSocket: \textbf{globalny kanał listy pokoi} (bez uwierzytelniania)
oraz \textbf{kanał czatu pokoju} (uwierzytelniany tokenem JWT w parametrze zapytania). Oba są
obsługiwane przez dwa rejestry połączeń w pamięci procesu: \code{ConnectionManager} i
\code{RoomChatManager} (\code{backend/ws/managers.py}). Handlery znajdują się w
\code{backend/routers/chat.py} (\code{websocket_rooms}, \code{websocket_chat}), serializery w
\code{backend/serializers.py} (\code{get_rooms_payload}, \code{_members_payload},
\code{_serialize_event_state}, \code{_msg_dict}), weryfikacja podpisu w \code{backend/security.py}
(\code{_chat_signing_message}, \code{_verify_chat_signature}), a stałe (\code{CHAT_SIGNATURE_SKEW_SECONDS},
\code{SYSTEM_SENDER_ADDRESS}, \code{WS_LIMITS}, \code{MAX_MSGS}, \code{WINDOW}) w
\code{backend/config.py}. Po stronie klienta odpowiadają im \code{roomsStore.ts}, \code{chatStore.ts}
i \code{signing.ts} (szczegóły w części o froncie).

\section{Menedżery połączeń}
\label{sec:rt-managery}

\subsection{\texttt{ConnectionManager} --- globalna lista pokoi}

\begin{tabularx}{\textwidth}{lY}
\toprule
Element & Opis \\
\midrule
\code{active_connections} & \code{list[WebSocket]}: każde podłączone gniazdo \code{/ws/rooms}. \\
\code{connect(ws)} & Korutyna: akceptuje gniazdo i dodaje je do listy. \\
\code{disconnect(ws)} & Usuwa gniazdo z listy (bez efektu, gdy go nie ma). \\
\code{broadcast(data: str)} & Korutyna: wysyła tekst do \emph{każdego} połączenia; gniazda, które zgłaszają wyjątek przy wysyłce, są usuwane. \\
\bottomrule
\end{tabularx}

Nieosiągalne przeglądarki (zamknięta karta, utrata sieci) są usuwane leniwie przy najbliższym
rozgłoszeniu. Po stronie serwera nie ma heartbeatu.

\subsection{\texttt{RoomChatManager} --- czat pokoju}

\begin{tabularx}{\textwidth}{lY}
\toprule
Element & Opis \\
\midrule
\code{rooms} & \code{dict[str, dict[WebSocket, str]]}: \code{nazwa_pokoju -> { WebSocket -> adres_tożsamości }}. \\
\code{connect(room, ws, address)} & Akceptuje gniazdo i rejestruje je pod kluczem pokoju. \\
\code{disconnect(room, ws)} & Wyrejestrowuje jedno gniazdo; klucz pokoju jest usuwany, gdy zostanie pusty. \\
\code{broadcast(room, payload: str)} & Wysyła tekst do \emph{każdego} połączenia w pokoju; nieaktywne gniazda są usuwane. \\
\code{disconnect_user(room, address, code=4409)} & Zamyka każde gniazdo uwierzytelnione jako \code{address} (bez rozróżniania wielkości liter). \\
\bottomrule
\end{tabularx}

\section{Kanał \texttt{/ws/rooms} --- globalna lista pokoi}
\label{sec:rt-rooms}

\begin{lstlisting}
GET /ws/rooms
\end{lstlisting}

\subsection{Połączenie}

\begin{itemize}[nosep]
  \item \textbf{Brak uwierzytelniania.} Gniazdo jest akceptowane bezwarunkowo i rejestrowane w \code{ConnectionManager}.
  \item Serwer natychmiast wysyła jako pierwszą ramkę bieżącą listę aktywnych pokoi w JSON.
  \item Następnie gniazdo trwa w pętli \code{receive_text()}, cicho zużywając dane przychodzące (żadna wiadomość od klienta nie wywołuje akcji).
\end{itemize}

\subsection{Ramki wychodzące}

Kanał wysyła jeden typ ramki: \textbf{goły tablicowy JSON} z podsumowaniami pokoi (bez opakowującego
pola \code{type}). Element tablicy:

\begin{lstlisting}[language=Python]
{
  "name": "duel-arena",
  "game": "Quake III Arena",
  "lobby_location": "eu-central",
  "players_active": 2,
  "players_max": 4,
  "member_addresses": ["0xAbc...", "0xDef..."],
  "description": "Looking for a duel partner",
  "communicator_link": "https://discord.gg/abc123",
  "requirements": "Microphone required",
  "expires_at": "2026-06-19T12:00:00Z"
}
\end{lstlisting}

\begin{tabularx}{\textwidth}{llY}
\toprule
Pole & Typ & Opis \\
\midrule
\code{name} & \code{string} & Unikalna nazwa pokoju. \\
\code{game} & \code{string} & Nazwa kategorii gry. \\
\code{lobby_location} & \code{string} & Kod lokalizacji lobby (np. \code{eu-central}; pełna lista w rozdziale o REST). \\
\code{players_active} & \code{number} & Bieżąca liczba członków (\code{len(r.members)}). \\
\code{players_max} & \code{number} & Pojemność ustawiona przy tworzeniu. \\
\code{member_addresses} & \code{string[]} & Adresy tożsamości aktualnych członków (suma kontrolna). \\
\code{description} & \code{string | null} & Opis pokoju. \\
\code{communicator_link} & \code{string | null} & Zewnętrzny link do komunikatora. \\
\code{requirements} & \code{string | null} & Wymagania. \\
\code{expires_at} & \code{string} & ISO 8601 z sufiksem \code{Z}. \\
\bottomrule
\end{tabularx}

\subsection{Kiedy następuje rozgłoszenie}

Wywołanie \code{ConnectionManager.broadcast(get_rooms_payload(session))} następuje po każdej
mutacji zmieniającej widoczny zbiór pokoi:

\begin{tabularx}{\textwidth}{lY}
\toprule
Zdarzenie & Endpoint / źródło \\
\midrule
Pokój utworzony & \code{POST /rooms} \\
Dołączenie do pokoju & \code{POST /rooms/{name}/join} \\
Opuszczenie pokoju & \code{POST /rooms/{name}/leave} \\
Członek usunięty & \code{POST /rooms/{name}/members/{addr}/kick} \\
Pokój zamknięty (admin) & \code{DELETE /rooms/{name}} \\
Gra usunięta (admin) & \code{DELETE /games/{name}} (zamyka wszystkie pokoje tej gry) \\
Pokój wygasł (tło) & \code{cleanup_expired_rooms_task}, co \code{ROOM_CLEANUP_INTERVAL_SECONDS} (domyślnie 60 s) \\
Zmiana czasu wygaśnięcia pokoju & pośrednio, przy ustawieniu/zmianie wydarzenia (\code{PUT /rooms/{name}/event}) \\
\bottomrule
\end{tabularx}

\subsection{Usuwanie wygasłych pokoi}

Wygasłe pokoje (\code{Room.expires_at <= now}) są usuwane synchronicznie wewnątrz
\code{get_rooms_payload()} przy \emph{każdym} wywołaniu, więc rozgłoszenie nigdy nie zawiera
nieaktualnych wpisów. Osobne zadanie w tle \code{cleanup_expired_rooms_task} (interwał
\code{ROOM_CLEANUP_INTERVAL_SECONDS}) usuwa pokoje (wraz z wiadomościami, wydarzeniem i RSVP),
które wygasły między kolejnymi pobraniami listy, i rozgłasza odświeżoną listę.

\textbf{Uwaga:} zadanie czyszczące \emph{nie} wysyła ramek \code{room_closed} na kanał czatu --- odświeża
tylko listę globalną. Połączenia czatu z wygasłym pokojem pozostają otwarte, dopóki klient nie wyśle
wiadomości (wtedy kontrola członkostwa kończy połączenie kodem 4403) lub się nie rozłączy.

\section{Kanał \texttt{/ws/rooms/\{room\_name\}/chat} --- czat pokoju}
\label{sec:rt-chat}

\begin{lstlisting}
GET /ws/rooms/{room_name}/chat?token=<JWT>
\end{lstlisting}

\subsection{Połączenie i uwierzytelnianie}

\begin{enumerate}[nosep]
  \item \textbf{JWT w parametrze zapytania.} Parametr \code{token} niesie JWT HS256 uzyskany przy logowaniu. Przeglądarki nie mogą ustawiać własnych nagłówków WebSocket, więc token przekazuje się w URL.
  \item Serwer dekoduje JWT funkcją \code{_decode_jwt()}, by wydobyć adres tożsamości (claim \code{sub}). Niepowodzenie $\to$ zamknięcie kodem \textbf{4401}.
  \item Pokój musi istnieć w bazie. Brak $\to$ kod \textbf{4404}.
  \item Wywołujący musi być aktualnym członkiem pokoju. Brak $\to$ kod \textbf{4403}.
\end{enumerate}

Po przejściu wszystkich kontroli gniazdo jest rejestrowane w \code{RoomChatManager.connect()}.

\subsection{Kody zamknięcia}

\begin{tabularx}{\textwidth}{llY}
\toprule
Kod & Znaczenie & Kiedy wysyłany \\
\midrule
4401 & Unauthorized & JWT brakuje, wygasł lub jest nieprawidłowy. \\
4403 & Forbidden --- nie członek & Wywołujący nie jest na liście członków (sprawdzane przy połączeniu \emph{oraz} po każdej wiadomości przychodzącej). \\
4404 & Not Found & Pokój o danej nazwie nie istnieje. \\
4409 & Członkostwo zakończone & \code{disconnect_user} po usunięciu przez administratora; ponowne połączenie kończy się kodem 4403. \\
4429 & Rate limited & Przekroczono limit wiadomości dla adresu. \\
\bottomrule
\end{tabularx}

We froncie kody \textbf{4403} i \textbf{4409} ustawiają \code{kicked = true} i wyłączają ponowne łączenie;
pozostałe kody uruchamiają ponowne łączenie z wykładniczym wycofywaniem (backoff).

\subsection{Ponowna kontrola członkostwa}

Po każdej wiadomości przychodzącej serwer odpytuje bezpośrednio tabelę łączącą \code{RoomMember}
(nie pamięć podręczną relacji ORM):

\begin{lstlisting}[language=Python]
membership = session.exec(
    select(RoomMember).where(
        RoomMember.room_id == room.id,
        RoomMember.user_id == user.id,
    )
).first()
\end{lstlisting}

Jeśli wiersza członkostwa nie ma, gniazdo jest zamykane kodem \textbf{4403}. Dzięki temu
usunięty lub wyrzucony członek nie może dalej pisać, nawet jeśli jego gniazdo pozostało otwarte.

\subsection{Limit wiadomości na adres}

Dla każdego adresu tożsamości działa okno przesuwne:

\begin{tabularx}{\textwidth}{llllY}
\toprule
Zmienna & Zmienna env & Domyślnie & & Opis \\
\midrule
\code{MAX_MSGS} & \code{WS_RATE_LIMIT_MAX_MESSAGES} & 10 & & Maks. liczba wiadomości w oknie na adres. \\
\code{WINDOW} & \code{WS_RATE_LIMIT_TIME_WINDOW_SECONDS} & 10 & & Długość okna w sekundach. \\
\bottomrule
\end{tabularx}

Stan limitu przechowuje globalny słownik \code{WS_LIMITS: dict[str, list[float]]} (w pamięci,
\textbf{niewspółdzielony} między procesami roboczymi). Przy każdej wiadomości przychodzącej serwer
usuwa znaczniki starsze niż \code{WINDOW}, sprawdza, czy liczba $\geq$ \code{MAX_MSGS}, a po
przekroczeniu zamyka gniazdo kodem \textbf{4429}. Kontrola odbywa się przed parsowaniem JSON, więc
liczą się także wiadomości niepoprawne.

\subsection{Typy ramek wychodzących}
\label{sec:rt-ramki-wychodzace}

Każda ramka wychodząca to obiekt JSON z polem \code{type} (łańcuch) i polami zależnymi od typu.

\begin{tabularx}{\textwidth}{l>{\raggedright\arraybackslash}p{0.33\textwidth}Y}
\toprule
\code{type} & Kiedy & Zawartość \\
\midrule
\code{history} & Zaraz po połączeniu, przed pętlą odczytu & Ostatnie 50 wiadomości w kolejności chronologicznej. \\
\code{message} & Po zapisie nowej wiadomości (także systemowej) & Pojedyncza wiadomość, do wszystkich w pokoju. \\
\code{event_update} & Wydarzenie utworzone, zmienione, odwołane lub zmiana RSVP przy opuszczeniu/kicku & Pełny stan wydarzenia lub \code{null}. \\
\code{rsvp_update} & Członek ustawił swoje RSVP & Pojedynczy wpis RSVP. \\
\code{roster_update} & Dołączenie, opuszczenie lub kick & Pełna aktualna lista członków. \\
\code{room_closed} & Administrator zamknął pokój lub usunął grę & Nazwa pokoju. \\
\code{member_kicked} & Administrator usunął członka & Dane usuniętego i właściciela pokoju. \\
\code{error} & Odrzucona wiadomość przychodząca & Pole \code{detail}. \\
\bottomrule
\end{tabularx}

\paragraph{\texttt{history}.} Zawiera 50 ostatnich wiadomości w porządku chronologicznym.

\begin{lstlisting}[language=Python]
{
  "type": "history",
  "messages": [
    {
      "id": 1,
      "sender_address": "0xAbc...",
      "sender_username": "player1",
      "content": "gg",
      "created_at": "2026-06-18T20:00:00Z",
      "signature": "0x1234...",
      "sent_at": "2026-06-18T20:00:00.000Z",
      "verified": true
    }
  ]
}
\end{lstlisting}

\paragraph{\texttt{message}.} Rozgłaszana do każdego połączenia w pokoju po zapisaniu nowej wiadomości.

\begin{lstlisting}[language=Python]
{
  "type": "message",
  "message": {
    "id": 42,
    "sender_address": "0xAbc...",
    "sender_username": "player1",
    "content": "Hello!",
    "created_at": "2026-06-18T20:01:00Z",
    "signature": "0x1234...",
    "sent_at": "2026-06-18T20:01:00.000Z",
    "verified": true
  }
}
\end{lstlisting}

Pola obiektu wiadomości (\code{_msg_dict}):

\begin{tabularx}{\textwidth}{llY}
\toprule
Pole & Typ & Opis \\
\midrule
\code{id} & \code{number} & Klucz główny (autoinkrementacja). \\
\code{sender_address} & \code{string} & Adres tożsamości nadawcy (suma kontrolna). \\
\code{sender_username} & \code{string} & Nazwa nadawcy (ustalana w chwili rozgłoszenia). \\
\code{content} & \code{string} & Treść wiadomości. \\
\code{created_at} & \code{string} & Czas serwera, ISO 8601 z \code{Z}. \\
\code{signature} & \code{string | null} & Podpis EIP-191 w hex lub \code{null} (wiadomość niepodpisana). \\
\code{sent_at} & \code{string | null} & Podpisany znacznik czasu klienta (ISO 8601) lub \code{null}. \\
\code{verified} & \code{boolean} & \code{true}, gdy wiadomość ma podpis; podpis jest weryfikowany przy zapisie, a wiadomości z błędnym podpisem w ogóle nie trafiają do bazy. \\
\code{kind} & \code{string} & \emph{(opcjonalne, tylko dla nadawcy systemowego)} \code{"system"}. \\
\bottomrule
\end{tabularx}

Wiadomość z \code{kind: "system"} jest wstrzykiwana przez serwer jako ogłoszenie (np. o usunięciu
członka przez administratora); jej \code{sender_address} to \code{"__system__"}, a nazwa nadawcy
to \code{"System"}.

\paragraph{\texttt{event\_update}.} Wysyłana przy utworzeniu, zmianie lub odwołaniu wydarzenia, a także
gdy zmienia się zbiór RSVP (opuszczenie pokoju lub kick członka z RSVP).

\begin{lstlisting}[language=Python]
{
  "type": "event_update",
  "event": {
    "starts_at": "2026-06-20T18:00:00Z",
    "ends_at": "2026-06-20T21:00:00Z",
    "created_by": "0xAbc...",
    "created_at": "2026-06-18T12:00:00Z",
    "updated_at": "2026-06-18T12:30:00Z",
    "rsvps": [
      { "address": "0xAbc...", "username": "player1",
        "status": "present", "updated_at": "2026-06-18T12:30:00Z" }
    ]
  }
}
\end{lstlisting}

Pole \code{event} to pełny stan wydarzenia albo \code{null} po odwołaniu. Pola obiektu: \code{starts_at},
\code{ends_at}, \code{created_by} (adres twórcy pokoju, który zaplanował wydarzenie), \code{created_at},
\code{updated_at} oraz \code{rsvps} (tablica wpisów, po jednym na członka, który odpowiedział).
Wpis RSVP: \code{address}, \code{username}, \code{status} (\code{"present"}, \code{"absent"} lub
\code{"maybe"}) i \code{updated_at}.

\paragraph{\texttt{rsvp\_update}.} Wysyłana, gdy pojedynczy członek ustawia RSVP
(\code{PUT /rooms/{name}/event/rsvp}).

\begin{lstlisting}[language=Python]
{
  "type": "rsvp_update",
  "rsvp": {
    "address": "0xDef...", "username": "player2",
    "status": "maybe", "updated_at": "2026-06-18T13:00:00Z"
  }
}
\end{lstlisting}

Frontend scala ramkę z lokalnym stanem wydarzenia, zastępując istniejący wpis z tego samego adresu
(deduplikacja adresów bez rozróżniania wielkości liter).

\paragraph{\texttt{roster\_update}.} Wysyłana przy dołączeniu, opuszczeniu lub usunięciu członka.

\begin{lstlisting}[language=Python]
{
  "type": "roster_update",
  "members": [
    { "address": "0xAbc...", "username": "player1", "is_admin": false },
    { "address": "0xDef...", "username": "player2", "is_admin": true }
  ]
}
\end{lstlisting}

Pole \code{members} to \emph{pełna} bieżąca lista (nie różnica) i zastępuje stan klienta. Element:
\code{address}, \code{username}, \code{is_admin} (czy adres jest na liście \code{ADMIN_ADDRESSES}).

\paragraph{\texttt{room\_closed}.} Wysyłana, gdy administrator usuwa pokój albo usuwa grę (co kaskadowo
zamyka wszystkie jej pokoje).

\begin{lstlisting}[language=Python]
{ "type": "room_closed", "room": "duel-arena" }
\end{lstlisting}

Frontend ustawia \code{closed = true}, zamyka gniazdo i wyłącza ponowne łączenie.

\paragraph{\texttt{member\_kicked}.} Wysyłana po usunięciu członka przez administratora
(\code{POST /rooms/{name}/members/{addr}/kick}).

\begin{lstlisting}[language=Python]
{
  "type": "member_kicked",
  "member_address": "0xAbc...",
  "member_username": "player1",
  "created_by": "0xAdmin...",
  "ownership_transferred": true
}
\end{lstlisting}

\begin{tabularx}{\textwidth}{llY}
\toprule
Pole & Typ & Opis \\
\midrule
\code{member_address} & \code{string} & Adres usuniętego członka. \\
\code{member_username} & \code{string} & Nazwa usuniętego członka. \\
\code{created_by} & \code{string} & Adres \emph{nowego} właściciela pokoju (po ewentualnym przekazaniu). \\
\code{ownership_transferred} & \code{boolean} & \code{true}, gdy usunięty był właścicielem pokoju. \\
\bottomrule
\end{tabularx}

Frontend aktualizuje \code{owner} na \code{created_by}; jeśli \code{member_address} zgadza się z adresem
bieżącej sesji, ustawia \code{kicked = true} i zamyka gniazdo. Bezpośrednio po tej ramce serwer wywołuje
\code{disconnect_user()} dla usuniętego adresu (kod zamknięcia 4409).

\paragraph{\texttt{error}.} Wysyłana, gdy walidacja po stronie serwera odrzuca wiadomość przychodzącą
(np. niepoprawny podpis).

\begin{lstlisting}[language=Python]
{ "type": "error", "detail": "signature_invalid" }
\end{lstlisting}

Frontend loguje ostrzeżenie; stan się nie zmienia. Odrzucona wiadomość \emph{nigdy} nie jest zapisywana
w bazie.

\subsection{Ramki przychodzące (klient $\to$ serwer)}
\label{sec:rt-ramki-przychodzace}

Klient wysyła \textbf{jeden} typ ramki --- wiadomość czatu: obiekt JSON z co najmniej polem \code{content}.

\textbf{Niepodpisana (starszy klient / klient bez klucza):}

\begin{lstlisting}[language=Python]
{ "content": "Hello, everyone!" }
\end{lstlisting}

\textbf{Podpisana (EIP-191):}

\begin{lstlisting}[language=Python]
{
  "content": "Hello, everyone!",
  "sent_at": "2026-06-18T20:01:00.000Z",
  "signature": "0x..."
}
\end{lstlisting}

\begin{tabularx}{\textwidth}{llY}
\toprule
Pole & Wymagane & Opis \\
\midrule
\code{content} & tak & Treść; przycinana, niepusta, $\leq$ 1000 znaków. \\
\code{sent_at} & z \code{signature} & Znacznik czasu klienta (ISO 8601) wchodzący do podpisywanego komunikatu. \\
\code{signature} & nie & Podpis EIP-191 w hex pod kanonicznym komunikatem (patrz następna sekcja). \\
\bottomrule
\end{tabularx}

\textbf{Potok walidacji po stronie serwera}
\begin{enumerate}[nosep]
  \item Kontrola członkostwa po każdej wiadomości --- brak członkostwa: zamknięcie kodem 4403.
  \item Kontrola limitu --- przekroczenie: zamknięcie kodem 4429.
  \item Parsowanie JSON; wiadomość niepoprawna jest ignorowana (bez odpowiedzi).
  \item Przycięcie \code{content}; pusta lub dłuższa niż 1000 znaków jest ignorowana (bez odpowiedzi i bez zamykania gniazda).
  \item Jeśli obecny jest \code{signature}: weryfikacja przez \code{_verify_chat_signature()}; przy niepowodzeniu serwer wysyła \code{{"type": "error", "detail": "signature_invalid"}} i pomija zapis (gniazdo \emph{nie} jest zamykane).
  \item Utworzenie wiersza \code{Message}, commit i rozgłoszenie \code{{"type": "message", "message": ...}} do wszystkich połączeń pokoju.
\end{enumerate}

\fig{sekwencja-czat.png}{0.85}{Sekwencja wysłania podpisanej wiadomości czatu (EIP-191): kontrola członkostwa i limitu, weryfikacja podpisu, zapis i rozgłoszenie.}

\section{Podpisywanie wiadomości (EIP-191)}
\label{sec:rt-podpis}

Wiadomości czatu mogą być podpisane kryptograficznie kluczem ECDSA nadawcy wyprowadzonym z BIP39,
dzięki czemu serwer dowodzi autorstwa end-to-end, a nie polega wyłącznie na JWT. Mechanizm jest
\textbf{opcjonalny}: klient bez klucza podpisującego wysyła wiadomości niepodpisane (\code{verified: false}).

\subsection{Kanoniczny komunikat}

Serwer i frontend budują dokładnie ten sam łańcuch:

\begin{lstlisting}
PlayLink signed chat message
room={room_name}
sent_at={sent_at}
content={content}
\end{lstlisting}

\begin{itemize}[nosep]
  \item \textbf{Białe znaki są istotne.} Cztery linie są łączone literalnym \texttt{\textbackslash n} (0x0A); bez końcowego znaku nowej linii.
  \item \code{room} to surowa nazwa pokoju (przy łączeniu kodowana w URL, przechowywana zdekodowana).
  \item \code{sent_at} to wygenerowany przez klienta łańcuch ISO 8601 przekazany w ładunku wiadomości.
  \item \code{content} to przycięta treść wiadomości.
\end{itemize}

Implementacja serwera (\code{_chat_signing_message}):

\begin{lstlisting}[language=Python]
def _chat_signing_message(room_name: str, content: str, sent_at: str) -> str:
    return (
        "PlayLink signed chat message\n"
        f"room={room_name}\n"
        f"sent_at={sent_at}\n"
        f"content={content}"
    )
\end{lstlisting}

Odpowiednik we froncie (\code{buildChatSigningMessage} w \code{signing.ts}):

\begin{lstlisting}[language=Java]
export function buildChatSigningMessage(room: string, content: string, sentAt: string): string {
    return `PlayLink signed chat message\nroom=${room}\nsent_at=${sentAt}\ncontent=${content}`;
}
\end{lstlisting}

Interfejs \code{ChatSigner} i użycie \code{signChatMessage} opisuje część o froncie.

\subsection{Weryfikacja podpisu}

Serwer weryfikuje podpis funkcją \code{_verify_chat_signature(room_name, content, sent_at, signature, address) -> bool}:

\begin{tabularx}{\textwidth}{lY}
\toprule
Krok & Szczegóły \\
\midrule
Parsowanie \code{sent_at} & Musi być poprawnym ISO 8601 (\code{datetime.fromisoformat()}). \\
Strefa czasowa & Znaczniki bez strefy są traktowane jako UTC. \\
Dryf zegara & \code{abs(now - sent_at)} musi być $\leq$ \textbf{120 sekund} (\code{CHAT_SIGNATURE_SKEW_SECONDS}) --- ochrona przed powtórzeniem (replay). \\
Odzyskanie sygnatariusza & \code{encode_defunct(text=canonical_message)}, następnie \code{Account.recover_message(message, signature)}. \\
Porównanie & Odzyskany adres (małe litery) porównywany z deklarowanym \code{address} (małe litery). \\
\bottomrule
\end{tabularx}

\subsection{Okno powtórzeń}

\code{CHAT_SIGNATURE_SKEW_SECONDS = 120} jest stałą w kodzie (nie jest konfigurowalna). Podpisy poza
oknem są odrzucane jako \code{signature_invalid}.

\subsection{\texttt{SYSTEM\_SENDER\_ADDRESS}}

Wiadomości generowane przez system (m.in. powiadomienia o usunięciu członka) używają stałego adresu
\code{"__system__"} i zawsze mają w ładunku \code{kind: "system"}. Są wstawiane przez serwer
i nie są podpisywane.

% =====================================================================

\chapter{Serwer MCP}
\label{ch:mcp}

Serwer Model Context Protocol (MCP) umożliwia asystentom AI (Claude Desktop, Claude Code, Cursor,
Antigravity) bezpieczną inspekcję platformy PlayLink w trybie \textbf{tylko do odczytu} (zgłoszenie \#95):
stanu platformy, obsługiwanych gier retro, aktywnych pokoi LFG i zdrowia regionów. Kod znajduje się
w \code{backend/mcp_server} (\code{server.py}, \code{__main__.py}). Serwer nie loguje się do PlayLink
i nie wykonuje żadnych mutacji --- odpytuje publiczne endpointy REST i, opcjonalnie, stos obserwowalności.
Miejsce MCP w planie realizacji i jego rola w ocenie projektu opisano w części o realizacji.

\fig{mcp.png}{0.9}{Serwer MCP uruchamiany lokalnie na komputerze użytkownika: odpytuje API lokalne (Docker) lub produkcyjne (Azure i Cloudflare) oraz opcjonalnie logi i metryki.}

\section{Narzędzia (tools)}
\label{sec:mcp-tools}

\begin{tabularx}{\textwidth}{p{5.8cm}Y}
\toprule
Narzędzie & Opis \\
\midrule
\code{list_games(region: str = "eu")} & Lista obsługiwanych gier retro. \\
\code{list_rooms(}\newline\code{region: str = "eu",}\newline\code{game: str | None = None)} & Aktywne pokoje LFG z opcjonalnym filtrem po grze. \\
\code{get_room_details(}\newline\code{room_name: str,}\newline\code{region: str = "eu")} & Szczegóły pokoju (gracze, komunikator, opis, wydarzenia). \\
\code{check_platform_health()} & Status endpointów \code{/healthz} i \code{/readyz} oraz opóźnienie dla wszystkich skonfigurowanych regionów. \\
\code{search_logs(}\newline\code{query: str = "level:ERROR",}\newline\code{minutes: int = 60,}\newline\code{limit: int = 50)} & Logi wszystkich regionów (VictoriaLogs, LogsQL), od najnowszych. \\
\code{query_metrics(query: str)} & Zapytanie PromQL do metryk (CPU, pamięć, zadania HTTP). \\
\bottomrule
\end{tabularx}

\section{Zasoby (resources)}
\label{sec:mcp-resources}

\subsection{Sparametryzowane per region}

\begin{tabularx}{\textwidth}{lY}
\toprule
URI & Opis \\
\midrule
\code{playlink://{region}/games} & Lista gier dla wybranego regionu (np. \code{playlink://eu/games}, \code{playlink://us/games}). \\
\code{playlink://{region}/rooms} & Aktywne pokoje w wybranym regionie (np. \code{playlink://eu/rooms}). \\
\code{playlink://{region}/rooms/{room_name}} & Szczegóły pokoju w wybranym regionie. \\
\bottomrule
\end{tabularx}

\subsection{Zbiorcze (multi-region)}

\begin{tabularx}{\textwidth}{lY}
\toprule
URI & Opis \\
\midrule
\code{playlink://games} & Lista gier ze sprawdzonych, dostępnych regionów. \\
\code{playlink://rooms} & Zagregowane podsumowanie aktywnych pokoi ze wszystkich skonfigurowanych regionów. \\
\code{playlink://rooms/{room_name}} & Wyszukiwanie pokoju we wszystkich skonfigurowanych regionach. \\
\code{playlink://status} & Stan zdrowia i dostępności bazy dla wszystkich regionów w JSON. \\
\bottomrule
\end{tabularx}

\section{Uruchamianie}
\label{sec:mcp-uruchomienie}

\textbf{stdio} (domyślny tryb dla asystentów AI):

\begin{lstlisting}[language=bash]
uv run --project backend python -m mcp_server
\end{lstlisting}

\textbf{sse / streamable-http} (sieciowo):

\begin{lstlisting}[language=bash]
uv run --project backend python backend/mcp_server/server.py --transport sse --port 8080
\end{lstlisting}

Parametry wiersza poleceń (\code{server.py}): \code{--transport} (\code{stdio}, \code{sse},
\code{streamable-http}; domyślnie \code{stdio}), \code{--host} (domyślnie \code{127.0.0.1}),
\code{--port} (domyślnie \code{8080}).

\section{Zmienne środowiskowe i konfiguracja regionów}
\label{sec:mcp-konfiguracja}

Regiony i adresy API są wykrywane dynamicznie, w następującej kolejności:

\begin{enumerate}[nosep]
  \item \code{PLAYLINK_REGIONS}: mapa JSON (np. \code{{"eu": "http://localhost:8000", "us": "http://localhost:8001"}}) albo lista rozdzielona przecinkami (\code{eu,us,asia}).
  \item Zmienne per region: \code{PLAYLINK_<REGION>_API} oraz \code{PUBLIC_REGION_<REGION>_API} (np. \code{PLAYLINK_EU_API}, \code{PUBLIC_REGION_US_API}).
  \item \code{PLAYLINK_DOMAIN}: po ustawieniu serwer wyprowadza adresy w formacie \code{https://playlink-<region>-api.<PLAYLINK_DOMAIN>}.
  \item Wartości domyślne: EU (\code{http://localhost:8000}), US (\code{http://localhost:8001}).
\end{enumerate}

\subsection{Logi i metryki (opcjonalnie)}

Narzędzia \code{search_logs} i \code{query_metrics} działają po ustawieniu \code{PLAYLINK_LOGS_URL},
\code{PLAYLINK_METRICS_URL}, \code{PLAYLINK_OBS_USER} (domyślnie \code{playlink}) oraz
\code{PLAYLINK_OBS_PASSWORD}; opis w części o obserwowalności (dokument
\code{docs/operations/observability.md}). Bez nich odpowiadają, że nie są skonfigurowane; pozostałe
narzędzia nie wymagają logowania. Na produkcji adresy mają postać
\code{https://playlink-logs.<domena>} i \code{https://playlink-metrics.<domena>} (basic auth, hasło
\code{OBS_PASSWORD} wdrożenia --- należy je uzyskać od opiekuna projektu i nigdy nie commitować).

\section{Konfiguracja w asystentach AI}
\label{sec:mcp-klienci}

Serwer MCP zawsze działa na komputerze użytkownika (Windows, macOS lub Linux) i komunikuje się
z API PlayLink wyłącznie przez HTTP. Różni się tylko miejsce, w którym działa API:

\begin{tabularx}{\textwidth}{llY}
\toprule
Przypadek & API & Potrzebne \\
\midrule
Lokalny & \code{localhost:8000} (EU) i \code{localhost:8001} (US) z \code{deploy/deploy.sh} & Docker \\
Zdalny & Wdrożone regiony, np. na Azure & Domena i opcjonalnie hasło do logów \\
\bottomrule
\end{tabularx}

\subsection{Instalacja narzędzi i kodu}

Potrzebne są Git oraz \code{uv} (menedżer pakietów Pythona, sam instaluje właściwą wersję Pythona):

\begin{tabularx}{\textwidth}{lYY}
\toprule
System & Git & uv \\
\midrule
Windows (PowerShell) & \code{winget install Git.Git} & \code{powershell -ExecutionPolicy ByPass -c "irm https://astral.sh/uv/install.ps1 | iex"} \\
macOS (Terminal) & \code{xcode-select --install} lub \code{brew install git} & \code{curl -LsSf https://astral.sh/uv/install.sh | sh} lub \code{brew install uv} \\
Linux (powłoka) & \code{sudo apt install git} (Debian/Ubuntu), \code{sudo dnf install git} (Fedora) & \code{curl -LsSf https://astral.sh/uv/install.sh | sh} \\
\bottomrule
\end{tabularx}

Po instalacji należy otworzyć nowy terminal i sprawdzić \code{git --version} oraz \code{uv --version}.
Następnie:

\begin{lstlisting}[language=bash]
git clone https://github.com/ioad-modul-sumatywny-26/playlink.git
cd playlink/backend
uv sync
pwd          # Windows PowerShell: (Get-Location).Path
\end{lstlisting}

Wypisana ścieżka to bezwzględna ścieżka do \code{backend}, używana dalej. Typowe wartości:

\begin{tabularx}{\textwidth}{lY}
\toprule
System & Przykład \\
\midrule
Windows & \code{C:/Users/you/playlink/backend} (ukośniki działają w JSON; przy odwrotnych ukośnikach należy je podwoić: \texttt{C:\textbackslash\textbackslash Users\textbackslash\textbackslash you\textbackslash\textbackslash...}) \\
macOS & \code{/Users/you/playlink/backend} \\
Linux & \code{/home/you/playlink/backend} \\
\bottomrule
\end{tabularx}

\subsection{Gdzie dodać serwer}

Wpisy JSON są takie same w każdym systemie; różni się tylko lokalizacja pliku i ścieżka do
\code{backend}. Blok \code{env} zmienia się między dwoma przypadkami poniżej.

\begin{tabularx}{\textwidth}{lY}
\toprule
Asystent & Gdzie dodać \\
\midrule
Claude Code & Jedno polecenie (patrz przypadki poniżej), takie samo we wszystkich systemach. \\
Claude Desktop & Settings $\to$ Developer $\to$ Edit Config. Windows: \texttt{\%APPDATA\%\textbackslash Claude\textbackslash claude\_desktop\_config.json}; macOS: \code{~/Library/Application Support/Claude/claude_desktop_config.json}. Brak oficjalnej aplikacji na Linuksa: użyć Claude Code lub Cursora. \\
Cursor & Settings $\to$ MCP lub \code{~/.cursor/mcp.json} (Windows: \texttt{\%USERPROFILE\%\textbackslash .cursor\textbackslash mcp.json}). \\
Antigravity / Gemini & \code{~/.gemini/config/mcp_config.json}. \\
\bottomrule
\end{tabularx}

Claude Desktop i Cursora należy po zapisaniu całkowicie zrestartować. Jeśli aplikacja nie znajduje
\code{uv} (typowe dla aplikacji desktopowych na macOS/Linux, które nie czytają profilu powłoki), w polu
\code{"command"} należy podać pełną ścieżkę: \code{which uv} (macOS/Linux) lub \code{where uv}
(Windows), np. \code{/Users/you/.local/bin/uv}.

\subsection{Przypadek 1: PlayLink uruchomiony lokalnie}

Najpierw należy uruchomić regiony (\code{deploy/deploy.sh eu}, opcjonalnie \code{deploy/deploy.sh us}).
Bez konfiguracji serwer używa \code{http://localhost:8000} (EU) i \code{http://localhost:8001} (US),
więc wpis nie potrzebuje \code{env}:

\begin{lstlisting}[language=bash]
{
  "mcpServers": {
    "playlink": {
      "command": "uv",
      "args": ["run", "--directory", "/path/to/playlink/backend", "python", "-m", "mcp_server"]
    }
  }
}
\end{lstlisting}

Claude Code:

\begin{lstlisting}[language=bash]
claude mcp add playlink -- uv run --directory /path/to/playlink/backend python -m mcp_server
\end{lstlisting}

Narzędzia można wypróbować bez asystenta, uruchamiając inspektora z katalogu głównego repozytorium
(Windows: PowerShell, te same polecenia):

\begin{lstlisting}[language=bash]
npx @modelcontextprotocol/inspector uv run --directory backend python -m mcp_server
npx @modelcontextprotocol/inspector --cli --method tools/list uv run --directory backend python -m mcp_server   # bez interfejsu przegladarkowego
\end{lstlisting}

Narzędzia \code{search_logs} i \code{query_metrics} działają lokalnie tylko wtedy, gdy uruchomiono także
stos obserwowalności; w przeciwnym razie odpowiadają, że nie są skonfigurowane.

\subsection{Przypadek 2: PlayLink uruchomiony zdalnie (produkcja)}

Dostęp do serwerów nie jest potrzebny. Serwer MCP wskazuje się na wdrożone regiony przez
\code{PLAYLINK_DOMAIN}; dla każdego regionu z \code{PLAYLINK_REGIONS} wyprowadza
\code{https://playlink-<region>-api.<domena>}:

\begin{lstlisting}[language=bash]
{
  "mcpServers": {
    "playlink": {
      "command": "uv",
      "args": ["run", "--directory", "/path/to/playlink/backend", "python", "-m", "mcp_server"],
      "env": {
        "PLAYLINK_DOMAIN": "<domain>",
        "PLAYLINK_REGIONS": "eu,us"
      }
    }
  }
}
\end{lstlisting}

Claude Code (Windows PowerShell, macOS i Linux jednakowo):

\begin{lstlisting}[language=bash]
claude mcp add playlink -e PLAYLINK_DOMAIN=<domain> -e PLAYLINK_REGIONS=eu,us -- uv run --directory /path/to/playlink/backend python -m mcp_server
\end{lstlisting}

Dla \code{search_logs} i \code{query_metrics} do \code{env} należy dodać poniższe zmienne (w Claude Code
po jednym \code{-e NAME=value} na zmienną):

\begin{lstlisting}[language=bash]
"PLAYLINK_LOGS_URL": "https://playlink-logs.<domain>",
"PLAYLINK_METRICS_URL": "https://playlink-metrics.<domain>",
"PLAYLINK_OBS_USER": "playlink",
"PLAYLINK_OBS_PASSWORD": "<OBS_PASSWORD>"
\end{lstlisting}

\subsection{Sprawdzenie połączenia}

Po restarcie asystenta należy poprosić go o uruchomienie \code{check_platform_health}; każdy
skonfigurowany region powinien zgłosić stan \code{reachable}.

\begin{tabularx}{\textwidth}{lY}
\toprule
Objaw & Prawdopodobna przyczyna \\
\midrule
Brak narzędzi PlayLink w asystencie & Błędna bezwzględna ścieżka w \code{--directory}, nieznaleziony \code{uv} (podać jego pełną ścieżkę) lub aplikacja nie została całkowicie zrestartowana. \\
Lokalnie: regiony nieosiągalne & Kontenery nie działają (\code{docker compose ls}). \\
\bottomrule
\end{tabularx}

\section{Bezpieczeństwo}
\label{sec:mcp-bezpieczenstwo}

\begin{itemize}[nosep]
  \item Serwer jest \textbf{tylko do odczytu}: wszystkie narzędzia i zasoby wykonują zapytania GET do publicznych endpointów API (\code{/games}, \code{/rooms}, \code{/rooms/{room_name}}, \code{/healthz}, \code{/readyz}) lub zapytania do logów i metryk; nie wymagają JWT i nie mają dostępu do kont ani do bazy danych.
  \item Domyślny transport \code{stdio} nie otwiera żadnego portu sieciowego; tryby sieciowe (\code{sse}, \code{streamable-http}) domyślnie wiążą się z \code{127.0.0.1}.
  \item Dostęp do logów i metryk wymaga hasła \code{PLAYLINK_OBS_PASSWORD} (basic auth, \code{OBS_PASSWORD} wdrożenia); należy je przechowywać poza repozytorium i przekazywać wyłącznie przez \code{env} konfiguracji klienta.
  \item Punkt \code{/metrics} backendu nie jest dostępny z internetu (żądania przez Cloudflare dostają \code{404}), więc metryki są dostępne tylko przez stos obserwowalności chroniony hasłem.
\end{itemize}

% SPRZECZNOSC: docs/backend/api-reference.md i realtime.md wskazują jako źródło backend/main.py
% (ConnectionManager, RoomChatManager, websocket_rooms, websocket_chat, _verify_chat_signature,
% stałe WS). Po refaktorze kod mieszka w routers/*.py, ws/managers.py, serializers.py, security.py
% i config.py; w tekście podano aktualne lokalizacje (zachowanie bez zmian).
% SPRZECZNOSC: api-reference.md pomija GET /metrics (router observability.py, poza schematem OpenAPI);
% dopisano go w rozdziale REST na podstawie kodu.
% SPRZECZNOSC: realtime.md podaje, że walidacja odrzuca „zniekształcony JSON" i pustą/zbyt długą
% treść; w kodzie routers/chat.py takie ramki są po cichu ignorowane (continue), bez ramki error i bez
% zamykania gniazda. Zastosowano zachowanie z kodu.
% SPRZECZNOSC: realtime.md opisuje pole verified jako „serwer odzyskał poprawnego sygnatariusza";
% w serializers._msg_dict verified = (signature is not None) --- weryfikacja następuje przy zapisie,
% niepoprawnie podpisane wiadomości nie są zapisywane.
% SPRZECZNOSC: realtime.md przykład lobby_location „us-west" nie istnieje na liście kodów
% (config.LOBBY_LOCATIONS: na-east, na-west, ...); użyto kodu eu-central.
% SPRZECZNOSC: README MCP w backend/mcp_server/README.md (przykład konfiguracji) nie zawiera zmiennych
% PLAYLINK_LOGS_URL itd. w JSON; pełny wariant (Przypadek 2) pochodzi z głównego README.md, nowszego.
% SPRZECZNOSC: sekcja „Bezpieczeństwo" (read-only, brak JWT, domyślny 127.0.0.1) została wywnioskowana
% z kodu server.py i opisu w README; backend/mcp_server/README.md nie ma osobnej sekcji o bezpieczeństwie.
